\documentclass[10pt]{article}
\usepackage[margin=0.76in]{geometry}
\usepackage{amsmath,amssymb,amsthm,mathtools}
\usepackage{graphicx}
\usepackage{xcolor}
\usepackage{microtype}
\usepackage[hidelinks]{hyperref}
\setlength{\parindent}{0pt}
\setlength{\parskip}{4.5pt}
\newtheorem{theorem}{Theorem}
\newtheorem{lemma}[theorem]{Lemma}
\newcommand{\R}{\mathbb R}
\newcommand{\F}{\mathcal F}
\newcommand{\C}{\mathcal C}
\newcommand{\relu}{\operatorname{ReLU}}
\newcommand{\step}{\operatorname{STEP}}

\title{Depth of the Minimal-Width Coding Memorizer\\
\large An exponential barrier, a precision-uniform optimality theorem, and a high-dimensional speedup for arXiv:2411.08735}
\author{Automated proof attempt\\
\texttt{agent\_lane\_03} \quad model: GPT5.6}
\date{11 August 2026}

\begin{document}
\maketitle

\begin{center}
\fcolorbox{blue!45!black}{blue!4}{\parbox{0.93\linewidth}{
\textbf{Status: substantial partial result, likely valid, pending human verification.}
For all finite-parameter piecewise-affine LReLU and stepped-LReLU classes in
the source, worst-case fixed-width depth is necessarily exponential in the
input coding parameter $K$.  If the depth cap is required to be independent
of output precision $M$, the source's $O(2^{d_xK})$ memorizer depth is
optimal up to width constants.  Conversely, a post-source ReLU memorization
theorem gives depth about the square root of that bound at the source's exact
minimal width whenever that width is at least $40$.  Low minimal widths and
networks allowed to exploit FLOOR remain open.}}
\end{center}

\section*{Source question}

The exact excerpt is Remark 79 and Table 3 on printed page 52 of the current
source paper (the following image is a direct crop of that page).

\begin{center}
\includegraphics[width=0.97\linewidth]{question_crop.png}
\end{center}

The source uses
\(C_n=\{j2^{-n}:0\le j<2^n\}\).  Its encoder bijects the
grid \(\C_K^{d_x}\) onto \(Z_K:=\C_{d_xK}\), so
\[
 N:=|Z_K|=2^{d_xK}.
\]
For a continuous target \(f:[0,1]^{d_x}\to[0,1]^{d_y}\), its memorizer
maps the corresponding point of \(Z_K\) to the encoded $M$-bit
quantization of $f$ at the grid point.  The source realizes this scalar
interpolation with depth $O(N)$, independently of $M$.

Call an activation family \emph{tame piecewise affine} if its graph is a
semialgebraic family with a bounded number of affine pieces and at most two
real activation parameters per unit.  This includes all of the source's
families
\[
 \F_+,\ \F_\pm,\ \F_{+,\mathfrak s},\ \F_{\pm,\mathfrak s},
 \ \F_{\alpha,\mathfrak s},\ \F_{\pm\alpha,\mathfrak s},
 \ \F^*_{+,1},\ \F^*_{\pm,1},
\]
and all their fixed-slope subfamilies.  It does not include FLOOR.

\begin{theorem}[Worst-case depth trichotomy]
\label{thm:main}
Fix scalar network width $w\ge1$, $d_x,d_y,K\ge1$, and a tame
piecewise-affine source activation family.  Let $N=2^{d_xK}$.

\begin{enumerate}
\item[(a)] If depth at most $L$ suffices to realize the exact coding
memorizer for every continuous target already at one output bit
($M=1$), then
\[
       L\ \ge\ c\,\frac{\sqrt N}{(w+1)^{3/2}}
       =c\,\frac{2^{d_xK/2}}{(w+1)^{3/2}},
\]
where $c>0$ is universal.  The same conclusion holds if exact matching is
replaced by uniform error strictly below $1/4$ on the coding grid.

\item[(b)] Suppose one depth cap $L=L(K)$, independent of $M$, realizes
the exact memorizer for every $M\ge1$ and every continuous target.  Then
\[
             L\ge \frac{N}{w(w+3)}.
\]
Thus the source's precision-uniform $O(N)$ memorizer is order-optimal for
fixed $w$.

\item[(c)] There is also a complementary upper bound.  For any fixed
\(\alpha\in(0,1)\), arbitrary coding data admit an exact
\(\sigma_\alpha\)-LReLU memorizer of width at most $40$ and depth
\[
 O\!\left(\sqrt{N B}+B\right),
 \qquad B:=1+d_xK+d_yM.
\]
Consequently, whenever the minimal width in the source is at least $40$,
this is an implementation \emph{at that minimal width}.  For
\(d_xK+d_yM=O(d_xK)\), parts (a) and (c) match in their exponential order
\(2^{d_xK/2}\), up to polynomial factors in $K$.
\end{enumerate}
\end{theorem}

\section*{Proof Intuition}

The target function is arbitrary, so its values on the finite input grid can
carry arbitrary labels.  Exact coding at one bit therefore forces the network
class to shatter all $N$ encoded grid points.  A width-$w$, depth-$L$
piecewise-affine network has only $O(w^2L)$ weights and $O(wL)$ nonlinear
units, so the sharp general VC upper bound is $O(w^3L^2)$.  This gives the
square-root exponential barrier.

Requiring the same depth bound for every $M$ is much stronger: dyadic label
vectors of all precisions are dense in the full $N$-cube.  The set of
evaluation vectors of bounded-depth tame networks is semialgebraic and has
dimension at most its $O(w^2L)$ real parameters.  Density forces dimension
$N$, hence a linear-in-$N$ depth bound.

The upper bound uses a March 2026 width--depth memorization theorem, then
converts every ReLU layer exactly into a fixed-LReLU layer at twice the width.

\section*{Proof}

Enumerate the input grid as $G_K=\{x_1,\ldots,x_N\}=\C_K^{d_x}$, and
write $z_j=\mathfrak e_K(x_j)\in Z_K$.  The $z_j$'s are distinct.
We first record why the target class really supplies arbitrary grid data.

\begin{lemma}[Arbitrary finite coding data]
\label{lem:data}
For every $M\ge1$ and every vector
\((a_1,\ldots,a_N)\in\C_M^N\), there is a continuous
\(f:[0,1]^{d_x}\to[0,1]^{d_y}\) whose coding memorizer satisfies
\[
       \mathfrak m_{K,M,f}(z_j)=a_j\qquad(1\le j\le N).
\]
\end{lemma}

\begin{proof}
Choose pairwise disjoint neighborhoods of the finitely many $x_j$'s and
continuous functions $0\le\eta_j\le1$, supported in those neighborhoods,
with \(\eta_j(x_j)=1\).  Set
\[
 f(x)=\left(\sum_{j=1}^N a_j\eta_j(x),0,\ldots,0\right).
\]
Disjointness of the supports keeps the first coordinate in $[0,1]$.
At $x_j$, its $M$-bit quantization is exactly $a_j$, and encoding a
vector whose only nonzero entry is its first entry leaves that entry
unchanged.  This is the asserted identity.
\end{proof}

For part (a), apply Lemma~\ref{lem:data} with labels
\(a_j\in\{0,1/2\}\subset\C_1\).  Thresholding the scalar network output at
$1/4$ realizes every binary labeling of $Z_K$, so $Z_K$ is shattered.

To count capacity, first eliminate the variable activation parameters.  For
all real \(\beta\),
\[
 \sigma_\beta(t)=\relu(t)-\beta\relu(-t),
 \qquad
 \rho_{\beta,s}(t)=\relu(t)-\beta\relu(-t)+s\step(t).
\]
Thus each source unit can be simulated by at most three units with fixed
piecewise-polynomial activations, with the affine layer after it performing
the displayed linear combination.  A width-$w$, depth-$L$ source network
is consequently contained in a fixed-activation network class with at most
\[
       W_*\le C_1(w+1)^2L \quad\hbox{parameters},\qquad
       U_*\le C_2wL \quad\hbox{nonlinear units}.
\]
Bartlett--Harvey--Liaw--Mehrabian's general piecewise-polynomial bound
\cite[Theorem 8]{BHLM} gives
\[
 \operatorname{VCdim}\le C_3W_*U_*\le C_4(w+1)^3L^2.
\]
Since the VC dimension is at least $N$, rearranging proves (a).  If the
grid error is less than $1/4$, thresholding still recovers the binary
labels, so the same proof applies.

For part (b), let $E_{w,L}\subset\R^N$ be the set of all evaluation vectors
\((\phi(z_1),\ldots,\phi(z_N))\) from widths at most $w$ and depths at
most $L$.  For each fixed depth and width profile, the graph of evaluation
in the affine weights, biases, and activation parameters is semialgebraic.
Tarski--Seidenberg therefore makes $E_{w,L}$ a finite union of
semialgebraic sets.  A depth-$n$ network has at most
\[
 n w(w+1)+2w(n-1)\le n w(w+3)
\]
real parameters.  Dimension cannot increase under a semialgebraic map, and
semialgebraic closure preserves dimension; hence
\[
       \dim\overline{E_{w,L}}\le Lw(w+3).
\]
The hypothesis and Lemma~\ref{lem:data} give
\(\C_M^N\subset E_{w,L}\) for every $M$.  The union of these dyadic grids
is dense in $[0,1]^N$, so
\([0,1]^N\subset\overline{E_{w,L}}\).  Its dimension is $N$, proving
\(N\le Lw(w+3)\) and part (b).

For part (c), apply Yang--Yang's Theorem 2.1 \cite{YY} to the $N$ scalar
inputs $z_j\in[0,1]$, whose separation is $1/N$, and to the
\(2^{d_yM}\) possible coded labels (rescale their integer indices at the
output).  In their explicit formula choose the bit-block parameter at least
both \(\log R\) and \(\log(2^{d_yM})\), where here
\(R=10N^3\sqrt\pi\), and choose the sample-block size as the ceiling of
the square-root balancing value.  Their displayed width is then at most
$20$, while their depth is
\[
       O\!\left(\sqrt{N(1+\log N+d_yM)}+1+\log N+d_yM\right).
\]
Since \(\log N=d_xK\log2\), this is the asserted bound with $B$.

Finally, a width-$r$ ReLU layer has an exact width-$2r$ realization using
one fixed positive LReLU, because
\[
       \relu(t)=\frac{\sigma_\alpha(t)+
       \alpha\sigma_\alpha(-t)}{1-\alpha^2}.
\]
At each layer the affine map produces $z$ and $-z$ in parallel, and the
next affine map performs the displayed combination.  This doubles width,
does not increase depth, and yields width $40$.  The activation
\(\sigma_\alpha\) belongs to every relevant variable positive/stepped
source family (take step size zero), completing (c). \qed

\section*{What is closed, and what remains open}

For tame LReLU families, part (a) gives a definitive qualitative answer to
the source's final suggestion that exponential scaling ``may be
unavoidable'': no worst-case implementation with depth subexponential in
$K$ exists at any fixed width.  Part (b) is sharp for the source's stronger,
implicit precision-uniform memorizer specification.  Part (c) shows that the
linear-in-$N$ construction can nevertheless be improved at fixed precision,
and at exact minimal width in dimensions where that width is at least $40$.

Three boundaries prevent classification as a full answer.  First, the exact
fixed-precision complexity at minimal widths below $40$, especially scalar
width one or two, is not settled.  Second, FLOOR is not semialgebraic with a
bounded number of pieces, so neither capacity proof applies to classes that
may use it.  This is substantive: with FLOOR and identity, a width-two,
constant-depth network can extract the $i$-th bit of a real parameter on
geometrically spaced inputs, giving infinite VC dimension.  Third, the
decoder's exponential dependence on $M$ is not closed here.

The durable attempt note records eight focused upgrade attempts, including
the failed width-one serialization and FLOOR-grid bit-extraction routes.

\section*{Reproducibility and literature boundary}

The script \path{code/verify_depth_bounds.py} checks the exact
ReLU-to-LReLU identity over rational test data, the parameter-count inequality,
and the constant-depth FLOOR bit-extraction obstruction.  These checks do not
replace the VC-dimension or semialgebraic-dimension proofs.

The source was revised in December 2025.  Yang--Yang was submitted in March
2026 and therefore postdates the source; its theorem is used only for part
(c).  Bartlett et al. supplies the pre-existing capacity theorem used in
part (a).  The precision-uniform closure argument and the application to this
specific coding question are proved above.  No duplicate packet or ledger
entry was found in the run indexes.  This is a bounded novelty check, not an
exhaustive priority claim.

\begin{thebibliography}{9}
\bibitem{RCG}
D. Rochau, R. Chan, and H. Gottschalk,
\emph{New advances in universal approximation with neural networks of
minimal width}, arXiv:2411.08735, current version December 2025.
Included as \texttt{source\_paper.pdf}; the exact two-page source excerpt is
also included as \texttt{source\_excerpt.pdf}.

\bibitem{BHLM}
P. L. Bartlett, N. Harvey, C. Liaw, and A. Mehrabian,
\emph{Nearly-tight VC-dimension and pseudodimension bounds for piecewise
linear neural networks}, JMLR 20 (2019), no. 63, 1--17;
arXiv:1703.02930.  Included under \texttt{references/}.

\bibitem{YY}
X. Yang and Y. Yang,
\emph{Memorization capacity of deep ReLU neural networks characterized by
width and depth}, arXiv:2603.09589 (10 March 2026).
Included under \texttt{references/}.
\end{thebibliography}

\end{document}
